% Fig 2.6 / spec Fig A02 -- pmf recovered from the jump heights of the cdf
% Source: mathstatch2.tex lines 320-349
\documentclass[tikz,border=2pt]{standalone}
\usepackage{amsmath}
\usepackage{amssymb}
\usepackage{xcolor}
\usetikzlibrary{arrows.meta}
\newcommand{\sssig}{\scriptscriptstyle}

\definecolor{journalink}{HTML}{1F1C18}
\definecolor{journalmuted}{HTML}{8A7F72}
\definecolor{journalaccent}{HTML}{9C2D2D}

\begin{document}
\begin{tikzpicture}[
  x=1cm,
  y=1cm,
  every node/.style={font=\normalsize, journalink},
  axis/.style={journalink, line width=0.8pt, -{Stealth[length=4pt]}},
  tick/.style={journalink, line width=0.8pt},
  stem/.style={journalmuted, line width=0.5pt},
  solidpt/.style={circle, fill=journalink, inner sep=0pt, minimum size=3.6pt}
]
  % ---- axes ----
  \draw[axis] (0,0) -- (6.5,0);
  \draw[axis] (0,0) -- (0,3.5);

  % ---- x ticks ----
  \foreach \x in {0,1,2,3,4,5,6}{
    \draw[tick] (\x,0.035) -- (\x,-0.106);
    \node[anchor=north] at (\x,-0.25) {$\x$};
  }

  % ---- y ticks (y=6 corresponds to probability 1) ----
  \draw[tick] (0.035,0) -- (-0.106,0) node[anchor=east, inner sep=2pt] {$0$};
  \draw[tick] (0.035,1) -- (-0.106,1) node[anchor=east, inner sep=2pt] {$1/6$};
  % Sixths are left unreduced so that the ticks match Fig. 2.1 tick for tick and
  % agree with the worked solution in the text, which writes p_{\sssig X}(3) = 2/6.
  % (were 1/3 and 1/2.)
  \draw[tick] (0.035,2) -- (-0.106,2) node[anchor=east, inner sep=2pt] {$2/6$};
  \draw[tick] (0.035,3) -- (-0.106,3) node[anchor=east, inner sep=2pt] {$3/6$};

  % ---- thin vertical stems (added for the web figure) ----
  \foreach \x/\y in {1/1, 2/1, 3/2, 4/1, 5/1}
    \draw[stem] (\x,0) -- (\x,\y);

  % ---- mass points ----
  \foreach \x/\y in {1/1, 2/1, 3/2, 4/1, 5/1}
    \node[solidpt] at (\x,\y) {};

  % ---- axis labels ----
  \node[anchor=west, inner sep=0pt] at (6.6,0) {$x$};
  \node[anchor=south, inner sep=0pt] at (0,3.6) {$p_{\sssig X}(x)$};
\end{tikzpicture}
\end{document}
